\section{Implementation and Evaluation}
\label{sec:benchmark}

For an empirical evaluation, we implemented the two masked-lookup modulus-conversion protocols by realizing $\func_\abb$ for a dishonest majority, using the SPDZ2k protocol~\cite{C:CDESX18}, i.e., authenticated additive secret sharing.
This places our numbers under the same security model as Fhenix~\cite{CCS:ZZLP25}, which we compare against and which is likewise dishonest-majority.
Following Fhenix, we implemented the offline and online parts separately:
\begin{enumerate}
    \item
        \emph{Offline phase.}
        Implemented in MP-SPDZ~\cite{EPRINT:Keller20}.
        In this phase all correlated randomness specific to a conversion is prepared: the mask together with the correction table, by \cref{protocol:modconv1-preproc} and \cref{protocol:modconv2-preproc} respectively.
    \item
        \emph{Online phase.}
        Also in MP-SPDZ, so that both phases are measured by the same instrumentation and over the same wire protocol.
        Mirroring Fhenix, whose precomputed gates are loaded into memory before timing starts, the preprocessed values are handed to the parties before the timer starts and that cost is excluded; the timer covers the conversion itself and nothing else.
\end{enumerate}

We convert from $q = 2^{63}$ to $t = 2^{64}$, the smallest nontrivial 2-power lift, chosen so that $q \mid t$ and \cref{remark:modconv1-divisible,remark:modconv2-divisible} both apply: $\shared{m}_q$ is never stored, following from any one of the corrections by a $\DivModConv$ command at no cost.
Programs are compiled for the ring $\ints_{2^{64}}$, so the machine ring is exactly $t$ and no slack bits are needed anywhere.
We use the same party counts $N \in \{4, 8, 16\}$ as Fhenix for the online phase, and $N = 2$ for the preprocessing, matching the setting in which \cite[Table~5]{CCS:ZZLP25} reports its own preprocessing costs.
Since $q$ and $\beta$ are 2-powers, the preprocessing of \cref{protocol:modconv1-preproc} loses its comparison, its acceptance test and its rejection loop, as the text following that protocol records; it then consists of $63$ shared random bits and local linear combinations.
For $\Pi_\mathsf{ModConv2}$, whose $\beta$ must satisfy $\beta \ge N+1$, we take the smallest 2-power above that floor, which keeps the precondition at $x \le q/2$ for every $N$, uniformly with $\Pi_\mathsf{ModConv1}$ at $\beta = 2$.

Every run checks its opened plaintext against the value computed in the clear, without a further opening.

\subsection{Decryption}
\label{sec:benchmark-online}

We measure the protocols as the decryption step they are meant to serve, not as a bare conversion.
Before the timer starts we sample a small secret key $\shared{\bm{s}} \in \ints_t^n$ and a public ciphertext $\bm{c} \in \ints_t^n$, and resample until $\langle \bm{c}, \bm{s}\rangle \bmod q$ meets the precondition of the conversion protocol;
accepting an instance is exactly the statement that the underlying error term lies in range.
We take $n = 1024$, matching the \gls{lwe} parameters $(n, q, p) = (1024, 2^{64}, 2)$ of \cite{CCS:ZZLP25}; as there, this is plain \gls{lwe} over $\ints_q$ rather than ring-\gls{lwe}, the secret key being a vector over $\ints_q$ and the first step of decryption an inner product.
The timed part is then
\[
    \shared{x}_t \coloneq \langle \bm{c}, \shared{\bm{s}} \rangle,
    \qquad
    \shared{x \bmod q}_q \gets \DivModConv(\shared{x}_t),
    \qquad
    \shared{x \bmod q}_t \gets \ModConv(\shared{x \bmod q}_q),
\]
followed by opening $\shared{m}_t \coloneq \shared{x}_t - \shared{x \bmod q}_t$, which yields $m = q \cdot \mathrm{bit}_{63}(x)$, the one-bit plaintext.
The $\DivModConv$ is free and is not even an instruction: the machine ring is $t$ and $q \mid t$, so the reduction is the identity on the sharing.

This costs two openings, one more than the conversion alone, and that extra one is deliberate.
The inner product and the final subtraction are local operations, but under SPDZ2k they still act on the MAC, and that cost is only realized when the run opens something;
timing the conversion in isolation would leave it out.
The inner product itself turns out to cost no communication whatever the dimension --- we measure the same five rounds and $0.36$ kB per opening at $n = 1$, $256$ and $1024$ --- since a linear combination and its MAC update are both local, so $n$ buys local computation only.
That local computation is not negligible, however: holding everything else fixed, the timed phase takes $0.014$~s at $n = 1$ against $0.037$~s at $n = 1024$, so the inner product is about $62\%$ of the online time.
It is therefore the dimension, and not the conversion, that governs how much a party's own cores matter to throughput.

\Cref{tab:benchmark-online} reports the result.
The measured round counts are MP-SPDZ's own, which count one round per polling iteration of its socket exchange rather than per network round-trip; the protocol-level figure is the two openings, against the three sequential openings of Fhenix's online phase.

\begin{table}[tbp]
    \caption{%
        Cost of one decryption, timed phase only, converting from $\ints_{2^{63}}$ to $\ints_{2^{64}}$.
        Openings are the protocol-level round complexity; rounds and data are measured by MP-SPDZ over the timed phase, under an emulated $1$~ms ping time and $1$~Gbit/sec bandwidth.%
    }
    \label{tab:benchmark-online}
    \centering
    \small
    \setlength{\tabcolsep}{4pt}
    \begin{tabular}{@{}llccccc@{}}
        \toprule
        Protocol & $N$ & $\beta$ & Openings & Rounds & Data/party & Time (ms) \\
        \midrule
        \multirow{3}{*}{$\Pi_\mathsf{ModConv1}$}
            & $4$  & $2$ & $2$ & $10$ & $0.72$ kB & $6.92$ \\
            & $8$  & $2$ & $2$ & $10$ & $1.68$ kB & $11.54$ \\
            & $16$ & $2$ & $2$ & $10$ & $3.60$ kB & $28.80$ \\
        \midrule
        \multirow{3}{*}{$\Pi_\mathsf{ModConv2}$}
            & $4$  & $8$  & $2$ & $10$ & $0.72$ kB & $6.87$ \\
            & $8$  & $16$ & $2$ & $10$ & $1.68$ kB & $8.42$ \\
            & $16$ & $32$ & $2$ & $10$ & $3.60$ kB & $18.99$ \\
        \bottomrule
    \end{tabular}
\end{table}

\subsection{Preprocessing}
\label{sec:benchmark-preproc}

\Cref{tab:benchmark-preproc} reports the cost of preparing the material for one conversion, at $N = 2$ and under an emulated $1$~ms ping time and $1$~Gbit/sec bandwidth, the conditions of \cite[Table~5]{CCS:ZZLP25}.

The preprocessing is where the two protocols separate.
$\Pi_\mathsf{ModConv1}$ issues no secure multiplication at all, but it needs $63$ shared random bits, and it is those bits that dominate its cost entirely.
$\Pi_\mathsf{ModConv2}$ instead has each party contribute its own mask, so it consumes no shared randomness and only the $(N-1)\beta^2 = 16$ multiplications of its convolution, and comes out an order of magnitude cheaper on both counts.
That comparison is not like for like, however, and \cref{sec:benchmark-caveats} says why.
It is worth noting where $\Pi_\mathsf{ModConv1}$'s cost actually sits: a shared random bit is not free in SPDZ2k, which realizes it by the square-root trick, drawing on Beaver triples generated over an auxiliary ring.
The protocol issues no $\Mul$ command, as \cref{tab:modconv} records, but the $\RandomBit$ commands it does issue are what the measurement is made of, so that is where an optimization of this preprocessing would have to bite.

\begin{table}[tbp]
    \caption{%
        Cost of the preprocessing for one modulus conversion, $N = 2$, under an emulated $1$~ms ping time and $1$~Gbit/sec bandwidth.
        Multiplication counts are exact, read off the MP-SPDZ compiler; the remaining columns are measured.%
    }
    \label{tab:benchmark-preproc}
    \centering
    \small
    \begin{tabular}{@{}lcccccc@{}}
        \toprule
        Protocol & $\beta$ & Mults & Random bits & Rounds & Data/party & Time (ms) \\
        \midrule
        $\Pi_\mathsf{ModConv1}$ & $2$ & $0$  & $63$ & $62$ & $12.1$ MB & $531$ \\
        $\Pi_\mathsf{ModConv2}$ & $4$ & $16$ & $0$  & $38$ & $0.36$ MB & $32.4$ \\
        \bottomrule
    \end{tabular}
\end{table}

These are the costs of a single conversion run in isolation, and they do not scale linearly: independent conversions collapse into shared rounds, so a per-thousand figure of the kind \cite[Table~5]{CCS:ZZLP25} reports cannot be obtained by multiplying \cref{tab:benchmark-preproc} by $1000$.
Batching one thousand conversions into a single program instead gives $37.1$~s, $1274$~MB and $2706$ rounds for $\Pi_\mathsf{ModConv1}$, against $8.85$~s, $336$~MB and $627$ rounds for $\Pi_\mathsf{ModConv2}$.

It is worth being precise about what the first of those numbers measures, because it is not our protocol.
Generating $63\,000$ shared random bits in SPDZ2k and doing nothing else costs $1272.76$~MB on the same host; the complete preprocessing of $\Pi_\mathsf{ModConv1}$ for the same one thousand conversions costs $1274.02$~MB.
The protocol's own contribution --- the corrections, the one-hot vector and the linear combinations --- is thus $1.26$~MB, or one part in a thousand, and the remaining $99.9\%$ is the price SPDZ2k charges for an authenticated random bit, about $20.2$~kB each.
Any protocol requiring $63$ such bits per conversion would pay the same, and a bespoke implementation would not pay it at all.
The defensible cost of this preprocessing is therefore the count of primitives it issues --- $63$ $\RandomBit$ commands and no $\Mul$ for $\Pi_\mathsf{ModConv1}$, $(N-1)\beta^2$ $\Mul$ commands and no $\RandomBit$ for $\Pi_\mathsf{ModConv2}$ --- with the megabytes read as a property of the framework rather than of the protocol.

\subsection{Network Conditions}
\label{sec:benchmark-network}

Fhenix emulates its network with \texttt{tc}, crossing ping times $\{1, 10, 100\}\,\mathrm{ms}$ with bandwidths $\{100\,\mathrm{Mbit/s}, 1\,\mathrm{Gbit/s}\}$, and \cref{tab:benchmark-latency} reports our figures over that same grid at $1$~Gbit/sec.
All parties run on one host and communicate over loopback, as Fhenix's own harness does; the emulated delay is therefore set to half the nominal ping time, since a loopback packet crosses the interface once in each direction.
Every run is preceded by a check that \texttt{netem} is in fact attached and by a measurement of the resulting loopback round-trip time, and the harness refuses to run unshaped rather than produce rows carrying a ping label they did not experience.
That guard exists because an earlier sweep did produce such rows: the emulation silently failed to attach, and the near-identical times at nominal $1$, $10$ and $100$~ms were the only sign.
The columns of \cref{tab:benchmark-latency} scale with the ping time as they should, which is what the guard is there to ensure.

One asymmetry should be recorded, because it runs against us.
Fhenix's single-ciphertext latency measurements are taken with a delay but no bandwidth limit at all, whereas every figure we report is additionally capped at $1$~Gbit/sec on a loopback shared by all $N$ parties.
The comparison in \cref{sec:benchmark-comparison} is therefore conservative in their favour.

\begin{table}[tbp]
    \caption{%
        Time (in ms) of one decryption, online phase, under emulated network conditions at $1$~Gbit/sec, following \cite[Table~2]{CCS:ZZLP25}.%
    }
    \label{tab:benchmark-latency}
    \centering
    \small
    \setlength{\tabcolsep}{4pt}
    \begin{tabular}{@{}llccc@{}}
        \toprule
        Protocol & Ping time & $4$ Parties & $8$ Parties & $16$ Parties \\
        \midrule
        \multirow{3}{*}{$\Pi_\mathsf{ModConv1}$}
            & $1$ ms   & $6.92$   & $11.54$  & $28.80$ \\
            & $10$ ms  & $61.21$  & $62.61$  & $84.90$ \\
            & $100$ ms & $605.47$ & $606.38$ & $666.22$ \\
        \midrule
        \multirow{3}{*}{$\Pi_\mathsf{ModConv2}$}
            & $1$ ms   & $6.87$   & $8.42$   & $18.99$ \\
            & $10$ ms  & $61.14$  & $67.14$  & $143.92$ \\
            & $100$ ms & $601.05$ & $607.93$ & $636.09$ \\
        \bottomrule
    \end{tabular}
\end{table}


\subsection{Throughput and Comparison with Fhenix}
\label{sec:benchmark-comparison}

\Cref{tab:benchmark-online,tab:benchmark-latency} time one decryption at a time.
Throughput is a different question, and it is the one \cite[Table~1]{CCS:ZZLP25} answers, so we measure it the same way: a batch of one thousand decryptions at $N = 4$ under a $1$~ms ping time and $1$~Gbit/sec bandwidth.

Our batch is vectorized rather than looped, so a thousand decryptions issue the same two opening \emph{instructions} as one.
The consequence is that the round count is essentially independent of the batch size --- $11$ measured rounds at a batch of $1000$ against $10$ at a batch of $1$ --- and the communication is $96$ bytes per decryption at $N = 4$, being two openings of a $\ints_{2^{64}}$ share to each of the $N-1$ peers.
On that basis we measure $32\,289$ decryptions per second for $\Pi_\mathsf{ModConv1}$ and $35\,804$ for $\Pi_\mathsf{ModConv2}$.

For the comparison we recomputed Fhenix's figures from the raw measurement logs published with its implementation rather than quoting its tables, so that both sides are derived the same way.
At $N = 4$ under the same nominal conditions its implementation achieves $73\,578$ decryptions per second, falling to $29\,281$ at $N = 8$ and $9\,447$ at $N = 16$.
We are thus within a factor of about two of it on throughput while ahead of it in round complexity, which is the expected shape of the comparison: their implementation is written directly in Rust for this one task, whereas ours is a program compiled to a generic multi-party virtual machine, and the gap between those two is not a statement about the protocols.
The figures above use a single thread per party; since the inner product accounts for some $62\%$ of the online time and is purely local arithmetic, a host with cores to spare should narrow the gap further, and \cref{sec:benchmark-caveats} records that we have not established by how much.

Latency, unlike throughput, does admit a like-for-like comparison, because it is governed by round trips rather than by how fast a party computes.
Fhenix's single-ciphertext measurements give $55.5$, $55.5$ and $56.5$~ms at $N = 4$, $8$ and $16$ under a $10$~ms ping time, and $555.8$, $566.0$ and $811.8$~ms under a $100$~ms ping time, that is to say about $5.5$ round trips, largely flat in the number of parties.
\Cref{tab:benchmark-latency} gives $61.2$, $62.6$ and $84.9$~ms, and $605.5$, $606.4$ and $666.2$~ms, for $\Pi_\mathsf{ModConv1}$ under the same two conditions --- about six round trips, and likewise flat.
The two protocols are therefore at parity on latency to within ten percent, and, as \cref{sec:benchmark-network} notes, our figures carry a bandwidth cap that theirs do not.

\subsection{Caveats}
\label{sec:benchmark-caveats}

Two qualifications apply to the numbers above.

The first is that the preprocessing of $\Pi_\mathsf{ModConv2}$ takes each party's contribution through SPDZ2k's own authenticated $\Input$ rather than through the $\CheckedInput$ of \cref{step:modconv2-preproc:input,step:modconv2-preproc:range} and the consistency check of \cref{step:modconv2-preproc:cons-check};
only the hot-one check of \cref{step:modconv2-preproc:hot-one} is performed.
Its figures in \cref{tab:benchmark-preproc} therefore understate the protocol as specified, by an amount these measurements do not determine, and its advantage over $\Pi_\mathsf{ModConv1}$ should not be read as a like-for-like comparison.
The preprocessing of $\Pi_\mathsf{ModConv1}$ takes no input and is unaffected.

The second concerns how much of the online timing is really online.
MP-SPDZ generates the authenticated randomness its MAC checks consume lazily, in batches whose size is a run-time parameter; a timed region that opens more values than that parameter allows will exhaust the pool and regenerate it mid-run, and the regeneration --- offline work --- is then billed to the online timer.
Left at its default this inflated our measured communication from $96$ bytes per decryption to $6405$, a factor of $64$, without any visible symptom other than the byte count itself.
The figures reported here are taken with the pool sized above the batch and with an untimed warm-up pass that forces the generation to happen before the timer starts, and every batched row was checked against the $96$ bytes per decryption that two openings must cost.
A measurement that is online by construction rather than by this kind of care would use MP-SPDZ's \texttt{Fake-Offline} facility with preprocessing read from disk, which we have not done.

The third is that we report one thread per party throughout.
The inner product is local arithmetic and accounts for about $62\%$ of the online time, so it is exactly the part that additional cores would parallelise, and Fhenix's harness does run several jobs per party concurrently.
Our own attempt to measure the effect was made on a host with fewer cores than parties, where it could only lose, and we therefore report the single-threaded figures and make no claim about what a wider host would give.

\claude{%
Every number above is measured, and \cref{tab:benchmark-online,tab:benchmark-latency} are now filled from runs with \texttt{tc} confirmed attached (\texttt{net\_emulated=yes} on every row, and the ping columns scale as they should).
Four things are still worth doing before this goes out.
(i) The throughput figures in \cref{sec:benchmark-comparison} come from the run you reported rather than from a committed CSV; re-run the throughput phase after the pool-sizing fix and confirm $32\,289$ and $35\,804$, or replace them.
(ii) Sweep \texttt{THREADS} on your host. The retracted single-threaded claim is the weakest point in \cref{sec:benchmark-caveats}, and with cores to spare the throughput gap to Fhenix should close; if it does, both that caveat and \cref{sec:benchmark-comparison} need rewriting in your favour.
(iii) Neither machine is named. Fhenix names a \texttt{c7a.32xlarge}; this section should name its counterpart, and the throughput comparison means little without both, since one is a Rust binary and the other a virtual machine.
(iv) I recomputed Fhenix's numbers from the logs in \texttt{benchmarks/}, not from its tables. Check one cell against the published table before relying on it, in case their reported figures are aggregated differently from the way I aggregated them (total ciphertexts over the slowest concurrent job).
}
